\section{Limitations}\label{sec:limitations}

\paragraph{Self-recovery only.} Every benchmark cohort was generated by the SR model itself. The benchmark therefore
measures how well the calibrator recovers SR parameters from SR data; it says nothing about whether the model describes
a given organism. Model adequacy has to be judged on each data set (fit by eye, observed against expected deaths,
leverage).

\paragraph{Sharp cohorts.} When deaths are tightly clustered in age (coefficient of variation of the death age below
0.05), the time step must be very small and fits take hours. Coverage stays at nominal (closed $\eta$ intervals covered
the truth on \C{sharp-eta-cov-k} of \C{sharp-eta-cov-n} pure sharp cohorts), but the intervals are wide: median closed
width \C{sharp-eta-width} decades for $\eta$ and \C{sharp-eps-width} for $\eps$. $\eta$ has no upper end on
\C{sharp-eta-open} of \C{sharp-n-pure} such cohorts. Raw $\beta$ is closed on only \C{sharp-beta-closed}; its upper end is
otherwise set by the bound on $\kappa$, not by the data, and removal must be read from $\beta\rbar(\kappa)/\eta$.

\paragraph{Fine time resolution.} Recording resolutions below about $6\times10^{-4}$ of the lifespan scale were not
benchmarked. This excludes, for example, \emph{E.~coli} lineages recorded every few minutes.

\paragraph{Human life tables.} On large human cohorts the upper ends of the $\eta$, $\beta$ and $\eps$ intervals are often
set by the search box rather than by the data, and must be read as ``no upper limit''. The claim shapes can also differ
between processors (Section~\ref{sec:software}). The grid rule also under-refines large human life tables: their grid error contracts more slowly than the rule assumes,
and at the grid it chooses only \C{conv-n-h-erule-le02} of \C{conv-n-reg-human} measured human cohorts were within
\C{nrule-target} nats (Section~\ref{sec:val-numerics}). The homogeneous model does not pin $\eta L^2$ for human life tables;
a heterogeneous fit is not part of the released calibrator.

\paragraph{External hazard.} The variant with a constant external hazard (H) recovers $\eta$ and $\eps$ on data with
external deaths, but $h$ itself is a nuisance parameter: it is not truth-scored, it can stay at its starting value,
and on sharp cohorts with a hazard H's likelihood was up to \C{h-sharp-offset-max} nats off at the truth
(measured on the previous library version). H's coverage was not measured on the sharp and $n$-ladder cohorts of
the benchmark. $h$ is never reported as an estimate.

\paragraph{Windowed fits.} The windowed calibrator W (Section~\ref{sec:estimation}) was scored on synthetic cohorts and on
cohorts simulated at real data sets, all from library v6 (Section~\ref{sec:val-benchmark}); no window-only cohort with an
unknown truth was scored, and the \C{w-sharp-untested} sharp windowed cohorts of the library (cohorts with a narrow spread of
death ages) were not tested. Its upper ends are often not set by the data: they sit at the search box or were widened after the
profile walk stopped short, in most closed claims at human life tables, in a large share of the cohorts planted at real tables
below $n=\C{w-large-n}$, and in about half of the $\eta$ claims at SR parameter grids. They must be read as lower bounds on the
interval's end at every sample size. The warning that W gives is keyed to $n\ge\C{w-large-n}$, and below that the claim labels
are the signal. W's reported optimum can lie up to \C{m-chi2-half} nats below the best point the fit saw; against H's fit at its
smallest hazard it fell short by up to \C{w-opt-a1-max} nats. The pre-registered rule on this gap failed, and its consequence,
withholding W, was waived after the result (Section~\ref{sec:val-benchmark}). On a human-size cohort the optimum differed
between two computers, and the second computer's optimum lay inside the first computer's claims (its own claims were not
recorded); the claims, not the point estimate, should be reported. All gate fits ran on one CPU family (AMD EPYC). The drop-one
leverage check was not tested for detection on windows. H's window in \code{sraging} was compared with the frozen calibrator
on two windowed cohorts, where the fits were equal, but its claims were not scored again on the benchmark; H's coverage on the
window arm (Section~\ref{sec:val-benchmark}) is that of the frozen calibrator. The window is one per cohort: per-individual entry ages are refused. The earlier envelope fit W0 reports a
range over a coarse grid instead of a profile interval, and its coverage of the truth has not been measured. Results that rest
on W0 should be read as provisional until they are compared with W on the same data.

\paragraph{Pinned $\kappa$.} The recommended pinned value $\kappa_0=10^{-3}$ rests on earlier versions of the
benchmark library (v1 and v5). The pinned variants were not rerun on library v6, which has more cohorts, sharp cohorts
and a different $\kappa$ grid.

\paragraph{Heterogeneity.} The released calibrators assume one homogeneous population. Benchmark cohorts with a
spread of thresholds are misspecified for every variant, and their coverage is far below nominal
(Section~\ref{sec:val-benchmark}). A heterogeneous likelihood exists in the engine but is not yet validated for
calibration.

\paragraph{Combinations.} The identified combinations ($\sqrt\eta/\eps$, $\beta\rbar(\kappa)$) were scored against the
truth for version 1.0.1 only on \C{twin-pop-n} cohorts of the previous benchmark library, through a twin of the
calibrator that also reports their intervals (Table~\ref{tab:identified}, note c). On the v6 benchmark only their
point errors are measured; their coverage there will be added in an update.

\paragraph{Leverage check.} The drop-one check's false-flag rate and power were measured for J only, on a few dozen
cohorts of independent individuals simulated from the model (Section~\ref{sec:val-leverage}). They were not measured
for H, for windowed fits, or separately for cohorts in which removal is strongly identified. A cluster of similar early
deaths is not seen by a check that drops one individual, and the effect fits of Section~\ref{sec:effects} have no
drop-one check.



\paragraph{Group comparisons.} The naming rule of Section~\ref{sec:effects} has measured error rates only for
independent individuals and for one design with individuals housed in cages (Section~\ref{sec:val-effects}). Its
seventh condition, the per-group cost of the joint fit, was added after the studies above had been analysed; its
simulated bars have not yet been measured in any study, so where it is applied the conservative $\chi^2$ bound is used.
No simulation has yet calibrated the naming rule (conditions 2--7) for individuals housed together (mice or mosquitoes in cages,
flies in vials, worms on plates) in designs with few such units of unequal size per group, or many small groups of units. Until it has, labels from such
designs that rest on design-effect scaling may have a false-positive rate above nominal, as the fly study found
(Section~\ref{sec:val-effects}), and should be read as provisional. This calibration will be added in an update of this paper. Results from other designs carry
the uncertainty of that calibration.
